% Options for packages loaded elsewhere
\PassOptionsToPackage{unicode}{hyperref}
\PassOptionsToPackage{hyphens}{url}
%
\documentclass[
]{article}
\usepackage{amsmath,amssymb}
\usepackage{iftex}
\ifPDFTeX
  \usepackage[T1]{fontenc}
  \usepackage[utf8]{inputenc}
  \usepackage{textcomp} % provide euro and other symbols
\else % if luatex or xetex
  \usepackage{unicode-math} % this also loads fontspec
  \defaultfontfeatures{Scale=MatchLowercase}
  \defaultfontfeatures[\rmfamily]{Ligatures=TeX,Scale=1}
\fi
\usepackage{lmodern}
\ifPDFTeX\else
  % xetex/luatex font selection
\fi
% Use upquote if available, for straight quotes in verbatim environments
\IfFileExists{upquote.sty}{\usepackage{upquote}}{}
\IfFileExists{microtype.sty}{% use microtype if available
  \usepackage[]{microtype}
  \UseMicrotypeSet[protrusion]{basicmath} % disable protrusion for tt fonts
}{}
\makeatletter
\@ifundefined{KOMAClassName}{% if non-KOMA class
  \IfFileExists{parskip.sty}{%
    \usepackage{parskip}
  }{% else
    \setlength{\parindent}{0pt}
    \setlength{\parskip}{6pt plus 2pt minus 1pt}}
}{% if KOMA class
  \KOMAoptions{parskip=half}}
\makeatother
\usepackage{xcolor}
\setlength{\emergencystretch}{3em} % prevent overfull lines
\providecommand{\tightlist}{%
  \setlength{\itemsep}{0pt}\setlength{\parskip}{0pt}}
\setcounter{secnumdepth}{-\maxdimen} % remove section numbering
\ifLuaTeX
  \usepackage{selnolig}  % disable illegal ligatures
\fi
\IfFileExists{bookmark.sty}{\usepackage{bookmark}}{\usepackage{hyperref}}
\IfFileExists{xurl.sty}{\usepackage{xurl}}{} % add URL line breaks if available
\urlstyle{same}
\hypersetup{
  pdftitle={Research Paper (Draft): ACP Mission Control, Agent Orchestration, and Swarm Warfare},
  hidelinks,
  pdfcreator={LaTeX via pandoc}}

\title{Research Paper (Draft): ACP Mission Control, Agent Orchestration,
and Swarm Warfare}
\author{}
\date{2026-02-12}

\begin{document}
\maketitle

\begin{quote}
Internal draft. \textbf{Non-published}.

This is a full rearchitecture of the annex into a
whitepaper/research-paper structure with topic-grouped evidence and
ontology-driven logic.
\end{quote}

\hypertarget{abstract}{%
\subsection{Abstract}\label{abstract}}

This paper argues that Agent Control Plane (ACP) capability is a
warfighting imperative because the decisive variable is no longer model
quality alone, but governed delegation throughput under contested
conditions. Across military and commercial ecosystems, agent operations
are converging on the same structural need: Mission Control for
autonomous execution.

The evidence shows six compounding realities:

\begin{enumerate}
\def\labelenumi{\arabic{enumi}.}
\tightlist
\item
  multi-agent orchestration is becoming the default scaling pattern,
\item
  swarms and low-cost autonomy compress operational timelines,
\item
  C2 integration debt remains a binding constraint,
\item
  governance/accountability obligations are hardening,
\item
  adversarial pressure concentrates at protocol, context, and connector
  boundaries,
\item
  acquisition pathways are already scaling autonomy faster than
  traditional governance mechanisms.
\end{enumerate}

Conclusion: ACP is not a feature layer. It is command infrastructure for
delegated autonomy.

\hypertarget{research-objective-and-decision-linkage}{%
\subsection{1) Research objective and decision
linkage}\label{research-objective-and-decision-linkage}}

\hypertarget{primary-research-objective}{%
\subsubsection{Primary research
objective}\label{primary-research-objective}}

Build a coherent, evidence-grounded model of how agent orchestration and
swarms change modern conflict, and translate that into:

\begin{itemize}
\tightlist
\item
  ACP-RA design primitives,
\item
  DAD organizational authorities,
\item
  near-term implementation priorities.
\end{itemize}

\hypertarget{decision-linkage}{%
\subsubsection{Decision linkage}\label{decision-linkage}}

\begin{itemize}
\tightlist
\item
  \textbf{ACP-RA feed:} what must be engineered and enforced in the
  control plane.
\item
  \textbf{DAD memo feed:} what must be institutionally governed,
  mandated, and audited.
\end{itemize}

\hypertarget{method-and-ontological-logic}{%
\subsection{2) Method and ontological
logic}\label{method-and-ontological-logic}}

\hypertarget{evidence-corpus}{%
\subsubsection{2.1 Evidence corpus}\label{evidence-corpus}}

This paper integrates three evidence layers:

\begin{enumerate}
\def\labelenumi{\arabic{enumi}.}
\tightlist
\item
  \textbf{Public operational/strategic evidence} (labs, policy, defense
  reporting, think tanks, standards signals).
\item
  \textbf{Academic evidence} (multi-agent systems, MARL, swarms,
  orchestration, safety/governance, adversarial resilience).
\item
  \textbf{Programmatic implementation evidence} (CJADC2 constraints,
  Replicator timelines, C-UAS scaling, protocol governance shifts).
\end{enumerate}

\hypertarget{ontology-entity-relation-model}{%
\subsubsection{2.2 Ontology (entity-relation
model)}\label{ontology-entity-relation-model}}

To avoid narrative drift and ensure coherence, claims are organized as
typed entities and relations:

\begin{itemize}
\tightlist
\item
  \textbf{Entity types:} \texttt{Claim}, \texttt{Evidence},
  \texttt{Constraint}, \texttt{ControlPrimitive}, \texttt{OrgAuthority},
  \texttt{Risk}, \texttt{Metric}.
\item
  \textbf{Core relations:}

  \begin{itemize}
  \tightlist
  \item
    \texttt{Evidence\ -\textgreater{}\ supports\ -\textgreater{}\ Claim}
  \item
    \texttt{Constraint\ -\textgreater{}\ limits\ -\textgreater{}\ ControlPrimitive}
  \item
    \texttt{Risk\ -\textgreater{}\ mitigated\_by\ -\textgreater{}\ ControlPrimitive}
  \item
    \texttt{Claim\ -\textgreater{}\ implies\ -\textgreater{}\ OrgAuthority}
  \item
    \texttt{Metric\ -\textgreater{}\ evaluates\ -\textgreater{}\ ControlPrimitive}
  \end{itemize}
\end{itemize}

\hypertarget{neo4j-supported-coherence-step}{%
\subsubsection{2.3 Neo4j-supported coherence
step}\label{neo4j-supported-coherence-step}}

Used local ontology + Neo4j sync to maintain structure and traceability
while rearchitecting:

\begin{itemize}
\tightlist
\item
  project node: \texttt{proj\_acp\_warfighting\_research}
\item
  paper node: \texttt{doc\_annex\_v6}
\item
  topic cluster nodes:

  \begin{itemize}
  \tightlist
  \item
    \texttt{doc\_topic\_mission\_control}
  \item
    \texttt{doc\_topic\_swarms}
  \item
    \texttt{doc\_topic\_governance}
  \end{itemize}
\end{itemize}

This paper is therefore organized as ontology topic clusters rather than
chronology.

\hypertarget{topic-cluster-a-mission-control-and-c2-compression}{%
\subsection{3) Topic cluster A: Mission Control and C2
compression}\label{topic-cluster-a-mission-control-and-c2-compression}}

\hypertarget{core-claim-a1}{%
\subsubsection{Core claim A1}\label{core-claim-a1}}

Warfighting advantage increasingly depends on reducing
delegation-to-effect latency while preserving bounded control.

\hypertarget{key-evidence-threads}{%
\subsubsection{Key evidence threads}\label{key-evidence-threads}}

\begin{itemize}
\tightlist
\item
  GAO CJADC2 reporting: progress exists, but framework, measurement, and
  integration barriers remain substantial. {[}35{]}
\item
  Frontier-style enterprise agent platforms frame deployment/governance
  stack as the bottleneck, not model IQ alone. {[}1{]}
\item
  Multi-agent platform announcements now emphasize orchestration, role
  delegation, and runtime boundaries. {[}49{]}{[}50{]}{[}51{]}
\end{itemize}

\hypertarget{implication-for-acp-ra}{%
\subsubsection{Implication for ACP-RA}\label{implication-for-acp-ra}}

ACP must implement Mission Control functions as first-class runtime
behaviors:

\begin{itemize}
\tightlist
\item
  intent packets,
\item
  delegation scopes,
\item
  policy checkpoints,
\item
  effect verification,
\item
  exception and rollback pathways.
\end{itemize}

\hypertarget{topic-cluster-b-swarms-volume-autonomy-and-force-survivability}{%
\subsection{4) Topic cluster B: Swarms, volume autonomy, and force
survivability}\label{topic-cluster-b-swarms-volume-autonomy-and-force-survivability}}

\hypertarget{core-claim-b1}{%
\subsubsection{Core claim B1}\label{core-claim-b1}}

Scale changes everything: command quality degrades unless orchestration
and policy enforcement are built for saturation conditions.

\hypertarget{key-evidence-threads-1}{%
\subsubsection{Key evidence threads}\label{key-evidence-threads-1}}

\begin{itemize}
\tightlist
\item
  CNAS: cheap drone mass creates persistent threat; layered C-UAS and
  scaled defenses are necessary. {[}54{]}
\item
  CNA: PRC swarm research is explicitly connected to Taiwan-relevant
  scenarios and lessons from ongoing conflicts. {[}55{]}
\item
  Army sustainment evidence shows swarm concepts moving into logistics
  and support-area security, not only maneuver units. {[}53{]}
\item
  Replicator 1/2 confirms autonomy and counter-autonomy are moving
  through accelerated fielding paths. {[}47{]}{[}56{]}
\end{itemize}

\hypertarget{implication-for-acp-ra-1}{%
\subsubsection{Implication for ACP-RA}\label{implication-for-acp-ra-1}}

ACP cannot assume sparse event flow. It must be saturation-resilient:

\begin{itemize}
\tightlist
\item
  queue-aware prioritization,
\item
  consequence-tier gating,
\item
  budget-aware action authorization,
\item
  deterministic degraded-mode operation.
\end{itemize}

\hypertarget{topic-cluster-c-interoperability-protocols-and-orchestration-ecosystems}{%
\subsection{5) Topic cluster C: Interoperability, protocols, and
orchestration
ecosystems}\label{topic-cluster-c-interoperability-protocols-and-orchestration-ecosystems}}

\hypertarget{core-claim-c1}{%
\subsubsection{Core claim C1}\label{core-claim-c1}}

Interoperability standards increase both coordination power and attack
surface. Protocol governance becomes decisive.

\hypertarget{key-evidence-threads-2}{%
\subsubsection{Key evidence threads}\label{key-evidence-threads-2}}

\begin{itemize}
\tightlist
\item
  A2A and MCP trajectory indicates rapid protocolization of
  agent-to-agent and tool-context interactions.
  {[}5{]}{[}6{]}{[}51{]}{[}52{]}
\item
  Linux Foundation governance move for A2A suggests ecosystem-wide
  standardization momentum. {[}52{]}
\item
  Multi-vendor enterprise platform competition is converging on
  orchestration + identity + runtime controls. {[}49{]}{[}50{]}{[}51{]}
\end{itemize}

\hypertarget{implication-for-acp-ra-2}{%
\subsubsection{Implication for ACP-RA}\label{implication-for-acp-ra-2}}

ACP should enforce protocol-bound trust, not app-bound trust:

\begin{itemize}
\tightlist
\item
  signed capability cards,
\item
  attested agent identity,
\item
  constrained delegation chains,
\item
  anti-replay and provenance at message boundaries.
\end{itemize}

\hypertarget{topic-cluster-d-safety-legal-accountability-and-command-responsibility}{%
\subsection{6) Topic cluster D: Safety, legal accountability, and
command
responsibility}\label{topic-cluster-d-safety-legal-accountability-and-command-responsibility}}

\hypertarget{core-claim-d1}{%
\subsubsection{Core claim D1}\label{core-claim-d1}}

At machine speed, accountability is impossible without architectural
provenance.

\hypertarget{key-evidence-threads-3}{%
\subsubsection{Key evidence threads}\label{key-evidence-threads-3}}

\begin{itemize}
\tightlist
\item
  Responsible-use policy frameworks are becoming lifecycle-oriented, not
  principle-only. {[}40{]}{[}41{]}{[}42{]}{[}43{]}{[}44{]}
\item
  ICRC/SIPRI logic remains central: responsibility is retained by humans
  and institutions; it cannot transfer to systems. {[}45{]}{[}46{]}
\item
  Frontier risk artifacts emphasize deceptive/failed tool-use conditions
  as practical risk domains. {[}12{]}{[}13{]}{[}14{]}
\end{itemize}

\hypertarget{implication-for-acp-ra-3}{%
\subsubsection{Implication for ACP-RA}\label{implication-for-acp-ra-3}}

Every high-consequence action path requires attributable evidence:

\begin{itemize}
\tightlist
\item
  who delegated,
\item
  what scope was active,
\item
  what policy checks fired,
\item
  what tools executed,
\item
  what verified effect occurred.
\end{itemize}

\hypertarget{topic-cluster-e-adversarial-resilience-and-contested-operation}{%
\subsection{7) Topic cluster E: Adversarial resilience and contested
operation}\label{topic-cluster-e-adversarial-resilience-and-contested-operation}}

\hypertarget{core-claim-e1}{%
\subsubsection{Core claim E1}\label{core-claim-e1}}

The highest-value attack surfaces are no longer only endpoints; they are
orchestration pathways and context channels.

\hypertarget{key-evidence-threads-4}{%
\subsubsection{Key evidence threads}\label{key-evidence-threads-4}}

\begin{itemize}
\tightlist
\item
  OpenClaw/ClawHub incidents and follow-on warnings demonstrate agent
  supply-chain attack behavior in the wild.
  {[}25{]}{[}26{]}{[}27{]}{[}28{]}
\item
  RAND and CSIS analyses stress adaptive adversaries, contested
  information advantage, and distributed mission command logic.
  {[}38{]}{[}39{]}
\item
  Adversarial MARL and deception-attack literature supports
  architecture-level robustness requirements.
\end{itemize}

\hypertarget{implication-for-acp-ra-4}{%
\subsubsection{Implication for ACP-RA}\label{implication-for-acp-ra-4}}

ACP should treat resilience as a control objective with explicit design
controls:

\begin{itemize}
\tightlist
\item
  partition-safe execution,
\item
  context provenance and confidence fusion,
\item
  connector attestation,
\item
  rapid revocation/reconstitution.
\end{itemize}

\hypertarget{topic-cluster-f-infrastructure-energy-and-industrial-base-constraints}{%
\subsection{8) Topic cluster F: Infrastructure, energy, and industrial
base
constraints}\label{topic-cluster-f-infrastructure-energy-and-industrial-base-constraints}}

\hypertarget{core-claim-f1}{%
\subsubsection{Core claim F1}\label{core-claim-f1}}

Compute and power constraints shape autonomy feasibility and tempo as
much as algorithmic capability.

\hypertarget{key-evidence-threads-5}{%
\subsubsection{Key evidence threads}\label{key-evidence-threads-5}}

\begin{itemize}
\tightlist
\item
  Capital and power are increasingly first-order autonomy
  enablers/constraints. {[}29{]}{[}30{]}{[}31{]}{[}32{]}
\item
  Operational concentration risk rises when autonomy depends on fragile
  infrastructure concentration.
\end{itemize}

\hypertarget{implication-for-acp-ra-5}{%
\subsubsection{Implication for ACP-RA}\label{implication-for-acp-ra-5}}

ACP must support heterogeneous model tiers and edge fallback profiles,
including low-connectivity and low-power execution envelopes.

\hypertarget{synthesis-acp-ra-design-primitives-ontology-consolidated}{%
\subsection{9) Synthesis: ACP-RA design primitives
(ontology-consolidated)}\label{synthesis-acp-ra-design-primitives-ontology-consolidated}}

The following primitives emerge as non-optional from cross-topic
synthesis.

\hypertarget{authority-primitives}{%
\subsubsection{9.1 Authority primitives}\label{authority-primitives}}

\begin{itemize}
\tightlist
\item
  Commander Intent Packet (CIP)
\item
  Delegation Authority Token (DAT)
\item
  Trust Scope Manifest (TSM)
\end{itemize}

\hypertarget{runtime-enforcement-primitives}{%
\subsubsection{9.2 Runtime enforcement
primitives}\label{runtime-enforcement-primitives}}

\begin{itemize}
\tightlist
\item
  action-time policy checks
\item
  consequence-tiered authorization
\item
  cross-agent interaction control
\end{itemize}

\hypertarget{evidence-primitives}{%
\subsubsection{9.3 Evidence primitives}\label{evidence-primitives}}

\begin{itemize}
\tightlist
\item
  action trajectory logging
\item
  evidence envelopes (context hash + tool receipts + effect
  verification)
\item
  forensic replay support
\end{itemize}

\hypertarget{resilience-primitives}{%
\subsubsection{9.4 Resilience primitives}\label{resilience-primitives}}

\begin{itemize}
\tightlist
\item
  degraded-mode profiles
\item
  partition-safe operation
\item
  rejoin and reconstitution controls
\end{itemize}

\hypertarget{supply-chain-primitives}{%
\subsubsection{9.5 Supply-chain
primitives}\label{supply-chain-primitives}}

\begin{itemize}
\tightlist
\item
  signed skills/connectors
\item
  provenance attestation
\item
  emergency revocation channels
\end{itemize}

\hypertarget{evaluation-primitives}{%
\subsubsection{9.6 Evaluation primitives}\label{evaluation-primitives}}

\begin{itemize}
\tightlist
\item
  mission-context benchmark suites
\item
  adversarial injection suites
\item
  drift-triggered rollback rules
\end{itemize}

\hypertarget{direct-feed-into-dad-memo-organizational-design}{%
\subsection{10) Direct feed into DAD memo (organizational
design)}\label{direct-feed-into-dad-memo-organizational-design}}

\hypertarget{why-dad-remains-structurally-necessary}{%
\subsubsection{10.1 Why DAD remains structurally
necessary}\label{why-dad-remains-structurally-necessary}}

No single functional lane can own delegated autonomy at force scale. DAD
is required as a cross-domain control authority integrating doctrine,
policy, and technical enforcement.

\hypertarget{minimum-dad-authorities-research-derived}{%
\subsubsection{10.2 Minimum DAD authorities
(research-derived)}\label{minimum-dad-authorities-research-derived}}

\begin{itemize}
\tightlist
\item
  approve/deny delegation scope baselines,
\item
  enforce no-conformance/no-scale rule,
\item
  mandate evidence and telemetry standards,
\item
  direct degraded-mode posture policy,
\item
  suspend tool registries or inter-agent channels,
\item
  run autonomy red-team and containment drills,
\item
  require pre-scale eval-pack conformance.
\end{itemize}

\hypertarget{initial-implementation-sequence}{%
\subsubsection{10.3 Initial implementation
sequence}\label{initial-implementation-sequence}}

\begin{itemize}
\tightlist
\item
  \textbf{Phase 1:} authority and baseline publication,
\item
  \textbf{Phase 2:} enforcement plumbing + controlled fieldings,
\item
  \textbf{Phase 3:} compulsory conformance and scored oversight.
\end{itemize}

\hypertarget{conclusions}{%
\subsection{11) Conclusions}\label{conclusions}}

\begin{enumerate}
\def\labelenumi{\arabic{enumi}.}
\tightlist
\item
  Multi-agent orchestration and swarm operations are not hypothetical
  future capabilities; they are current operational realities.
\item
  C2 and mission-control quality now bound autonomy effectiveness more
  than isolated model quality.
\item
  Protocol and connector surfaces are becoming primary trust boundaries.
\item
  Governance without runtime enforcement is inadequate under
  machine-speed delegation.
\item
  ACP-RA should be treated as warfighting infrastructure, and DAD as its
  institutional control authority.
\end{enumerate}

\hypertarget{open-source-intake-queue-auth-blocked-extraction}{%
\subsection{12) Open-source intake queue (auth-blocked
extraction)}\label{open-source-intake-queue-auth-blocked-extraction}}

\begin{itemize}
\tightlist
\item
  \url{https://x.com/vraserx/status/2021654971343610058?s=52}
\item
  \url{https://x.com/theinformation/status/2021707053727662407?s=46}
\end{itemize}

Status: captured as evidence pointers; full text extraction blocked in
this runtime by X authentication constraints.

\hypertarget{policyoperational-source-references}{%
\subsection{13) Policy/operational source
references}\label{policyoperational-source-references}}

\begin{enumerate}
\def\labelenumi{\arabic{enumi}.}
\tightlist
\item
  \textbf{OpenAI} --- Introducing OpenAI Frontier (Feb 2026)\\
  \url{https://openai.com/index/introducing-openai-frontier/}
\item
  \textbf{TechCrunch} --- OpenAI launches a way for enterprises to build
  and manage AI agents (Feb 2026)\\
  \url{https://techcrunch.com/2026/02/05/openai-launches-a-way-for-enterprises-to-build-and-manage-ai-agents/}
\item
  \textbf{Reuters} --- OpenAI launches Codex app to gain ground in AI
  coding race (Feb 2, 2026)\\
  \url{https://www.reuters.com/business/media-telecom/openai-launches-codex-app-gain-ground-ai-coding-race-2026-02-02/}
\item
  \textbf{VentureBeat} --- OpenAI launches a Codex desktop app for macOS
  to run multiple AI coding agents in parallel (Feb 2026)\\
  \url{https://venturebeat.com/orchestration/openai-launches-a-codex-desktop-app-for-macos-to-run-multiple-ai-coding/}
\item
  \textbf{Google Developers Blog} --- A2A: A new era of agent
  interoperability (Apr 9, 2025)\\
  \url{https://developers.googleblog.com/en/a2a-a-new-era-of-agent-interoperability/}
\item
  \textbf{Model Context Protocol Blog} --- First MCP anniversary / spec
  cadence update (Nov 25, 2025)\\
  \url{https://blog.modelcontextprotocol.io/posts/2025-11-25-first-mcp-anniversary/}
\item
  \textbf{Google Cloud Docs} --- Gemini CLI (tools, ReAct loop, MCP
  servers)\\
  \url{https://docs.cloud.google.com/gemini/docs/codeassist/gemini-cli}
\item
  \textbf{Google Developers Blog} --- Build with Google Antigravity (Nov
  20, 2025)\\
  \url{https://developers.googleblog.com/build-with-google-antigravity-our-new-agentic-development-platform/}
\item
  \textbf{Anthropic Research} --- Building Effective Agents (Dec 19,
  2024)\\
  \url{https://www.anthropic.com/research/building-effective-agents}
\item
  \textbf{Anthropic Engineering} --- Effective context engineering for
  AI agents (Sep 29, 2025)\\
  \url{https://www.anthropic.com/engineering/effective-context-engineering-for-ai-agents}
\item
  \textbf{Anthropic Engineering} --- Demystifying evals for AI agents
  (Jan 9, 2026)\\
  \url{https://www.anthropic.com/engineering/demystifying-evals-for-ai-agents}
\item
  \textbf{Anthropic} --- Sabotage Risk Report: Claude Opus 4.6 (PDF, Feb
  2026)\\
  \url{https://www-cdn.anthropic.com/f21d93f21602ead5cdbecb8c8e1c765759d9e232.pdf}
\item
  \textbf{Anthropic} --- Claude Opus 4.6 System Card (PDF, Feb 2026)\\
  \url{https://www-cdn.anthropic.com/0dd865075ad3132672ee0ab40b05a53f14cf5288.pdf}
\item
  \textbf{Anthropic} --- Responsible Scaling Policy updates (Feb 2026
  update context)\\
  \url{https://www.anthropic.com/rsp-updates}
\item
  \textbf{Dario Amodei} --- The Adolescence of Technology (Jan 2026)\\
  \url{https://www.darioamodei.com/essay/the-adolescence-of-technology}
\item
  \textbf{Financial Times} --- Humanity needs to wake up to dangers of
  AI (Jan 26, 2026)\\
  \url{https://www.ft.com/content/c3098552-7204-4a93-844c-1b8569c9dcb2}
\item
  \textbf{Associated Press} --- Former OpenAI leader says safety took a
  backseat (May 17, 2024)\\
  \url{https://apnews.com/article/openai-jan-leike-safety-ilya-8a7ba341e06a66e9a7935bb06214edcb}
\item
  \textbf{Reuters} --- OpenAI sets up safety and security committee (May
  28, 2024)\\
  \url{https://www.reuters.com/technology/openai-sets-up-safety-security-committee-2024-05-28/}
\item
  \textbf{Reuters} --- OpenAI technology chief Mira Murati to leave (Sep
  25, 2024)\\
  \url{https://www.reuters.com/technology/artificial-intelligence/openais-technology-chief-mira-murati-leave-2024-09-25/}
\item
  \textbf{Reuters} --- Sutskever safety startup SSI raises \$1B (Sep 4,
  2024)\\
  \url{https://www.reuters.com/technology/artificial-intelligence/openai-co-founder-sutskevers-new-safety-focused-ai-startup-ssi-raises-1-billion-2024-09-04/}
\item
  \textbf{Platformer} --- OpenAI mission alignment team report (Feb 11,
  2026)\\
  \url{https://www.platformer.news/openai-mission-alignment-team-joshua-achiam/}
\item
  \textbf{The Verge} --- OpenAI reportedly disbanded mission alignment
  team (Feb 11, 2026)\\
  \url{https://www.theverge.com/ai-artificial-intelligence/877208/openai-reportedly-disbanded-its-mission-alignment-team}
\item
  \textbf{TechCrunch} --- Microsoft CEO says up to 30\% of code written
  by AI (Apr 29, 2025)\\
  \url{https://techcrunch.com/2025/04/29/microsoft-ceo-says-up-to-30-of-the-companys-code-was-written-by-ai/}
\item
  \textbf{Business Insider} --- Microsoft CTO: 95\% of code AI-generated
  in five years (Apr 2025)\\
  \url{https://www.businessinsider.com/microsoft-cto-ai-generated-code-software-developer-job-change-2025-4}
\item
  \textbf{The Verge} --- OpenClaw skill extensions security nightmare
  (Feb 2026)\\
  \url{https://www.theverge.com/news/874011/openclaw-ai-skill-clawhub-extensions-security-nightmare}
\item
  \textbf{Reuters} --- China warns of security risks linked to OpenClaw
  (Feb 5, 2026)\\
  \url{https://www.reuters.com/world/china/china-warns-security-risks-linked-openclaw-open-source-ai-agent-2026-02-05/}
\item
  \textbf{1Password} --- From magic to malware: OpenClaw skills as
  attack surface (Feb 2, 2026)\\
  \url{https://1password.com/blog/from-magic-to-malware-how-openclaws-agent-skills-become-an-attack-surface}
\item
  \textbf{Trend Micro} --- What OpenClaw reveals about agentic
  assistants (Feb 6, 2026)\\
  \url{https://www.trendmicro.com/en_us/research/26/b/what-openclaw-reveals-about-agentic-assistants.html}
\item
  \textbf{Reuters} --- Deals showing AI runs on capital (Feb 6, 2026)\\
  \url{https://www.reuters.com/technology/spacex-nvidia-deals-showing-ai-runs-capital-2026-02-06/}
\item
  \textbf{International Energy Agency} --- Energy demand from AI\\
  \url{https://www.iea.org/reports/energy-and-ai/energy-demand-from-ai}
\item
  \textbf{Reuters} --- Meta begins construction of \$10B Indiana data
  center (Feb 11, 2026)\\
  \url{https://www.reuters.com/business/meta-begins-construction-10-billion-indiana-data-center-boost-ai-capabilities-2026-02-11/}
\item
  \textbf{Reuters} --- Anthropic to shoulder some data-center expansion
  costs (Feb 11, 2026)\\
  \url{https://www.reuters.com/technology/anthropic-shoulder-some-costs-data-center-expansions-threaten-raise-power-bills-2026-02-11/}
\item
  \textbf{Reuters} --- Artificial Intelligencer: AI and politics at
  Davos (Jan 22, 2026)\\
  \url{https://www.reuters.com/technology/artificial-intelligence/artificial-intelligencer-how-ai-politics-dominated-davos-2026-01-22/}
\item
  \textbf{The Verge} --- Anthropic turns to skills for workplace utility
  (Oct 16, 2025)\\
  \url{https://www.theverge.com/ai-artificial-intelligence/800868/anthropic-claude-skills-ai-agents}
\item
  \textbf{GAO} --- Defense Command and Control: Further Progress Hinges
  on Establishing a Comprehensive Framework (Apr 2025)\\
  \url{https://files.gao.gov/reports/GAO-25-106454/index.html}
\item
  \textbf{CSIS} --- Ukraine's Future Vision and Current Capabilities for
  Waging AI-Enabled Autonomous Warfare (Mar 2025)\\
  \url{https://www.csis.org/analysis/ukraines-future-vision-and-current-capabilities-waging-ai-enabled-autonomous-warfare}
\item
  \textbf{CSIS} --- How Russia Is Reshaping Command and Control for
  AI-Enabled Warfare (Feb 2026)\\
  \url{https://www.csis.org/analysis/how-russia-reshaping-command-and-control-ai-enabled-warfare}
\item
  \textbf{RAND} --- How Could Artificial Intelligence Shape the Future
  of War? (2026)\\
  \url{https://www.rand.org/pubs/research_reports/RRA4316-1.html}
\item
  \textbf{RAND} --- Strategic competition in the age of AI: Emerging
  risks and opportunities from military use of AI (2024)\\
  \url{https://www.rand.org/pubs/research_reports/RRA3295-1.html}
\item
  \textbf{NATO} --- Summary of NATO's revised AI strategy (Jul 2024)\\
  \url{https://www.nato.int/en/about-us/official-texts-and-resources/official-texts/2024/07/10/summary-of-natos-revised-artificial-intelligence-ai-strategy}
\item
  \textbf{UK MOD (GOV.UK)} --- JSP 936: Dependable Artificial
  Intelligence in defence (Part 1 directive) (Nov 2024)\\
  \url{https://www.gov.uk/government/publications/jsp-936-dependable-artificial-intelligence-ai-in-defence-part-1-directive}
\item
  \textbf{UK MOD (GOV.UK)} --- Laying the Groundwork: Responsible AI
  Senior Officers' Report 2025 (Oct 2025)\\
  \url{https://www.gov.uk/government/publications/laying-the-groundwork-responsible-ai-senior-officers-report-2025}
\item
  \textbf{U.S. Department of State} --- Political Declaration on
  Responsible Military Use of Artificial Intelligence and Autonomy\\
  \url{https://www.state.gov/bureau-of-arms-control-deterrence-and-stability/political-declaration-on-responsible-military-use-of-artificial-intelligence-and-autonomy}
\item
  \textbf{UNODA} --- Artificial intelligence in the military domain
  (incl.~UNGA 79/239 context)\\
  \url{https://disarmament.unoda.org/en/our-work/emerging-challenges/artificial-intelligence-military-domain}
\item
  \textbf{SIPRI} --- Retaining Human Responsibility in the Development
  and Use of Autonomous Weapon Systems (Oct 2022)\\
  \url{https://www.sipri.org/publications/2022/policy-reports/retaining-human-responsibility-development-and-use-autonomous-weapon-systems-accountability}
\item
  \textbf{ICRC Law \& Policy Blog} --- Three lessons on AWS regulation
  to ensure accountability for IHL violations (Mar 2023)\\
  \url{https://blogs.icrc.org/law-and-policy/2023/03/02/three-lessons-autonomous-weapons-systems-ihl/}
\item
  \textbf{DIU} --- The Replicator Initiative (updated 2025)\\
  \url{https://www.diu.mil/replicator}
\item
  \textbf{Arms Control Association} --- AI and Nuclear Command and
  Control: It's Even More Complicated Than You Think (Sep 2025)\\
  \url{https://www.armscontrol.org/act/2025-09/features/artificial-intelligence-and-nuclear-command-and-control-its-even-more}
\item
  \textbf{Microsoft Copilot Blog} --- Multi-agent orchestration
  announcements at Build 2025\\
  \url{https://www.microsoft.com/en-us/microsoft-copilot/blog/copilot-studio/multi-agent-orchestration-maker-controls-and-more-microsoft-copilot-studio-announcements-at-microsoft-build-2025/}
\item
  \textbf{Google Cloud Blog} --- Build and manage multi-system agents
  with Vertex AI (Apr 2025)\\
  \url{https://cloud.google.com/blog/products/ai-machine-learning/build-and-manage-multi-system-agents-with-vertex-ai}
\item
  \textbf{Google Cloud Blog} --- Agent2Agent protocol (A2A) is getting
  an upgrade (Jul 2025)\\
  \url{https://cloud.google.com/blog/products/ai-machine-learning/agent2agent-protocol-is-getting-an-upgrade}
\item
  \textbf{Google Developers Blog} --- Google Cloud donates A2A to Linux
  Foundation (Jun 2025)\\
  \url{https://developers.googleblog.com/en/google-cloud-donates-a2a-to-linux-foundation/}
\item
  \textbf{U.S. Army} --- Swarm Technology in Sustainment Operations (Jan
  2025)\\
  \url{https://www.army.mil/article/282467/swarm_technology_in_sustainment_operations}
\item
  \textbf{CNAS} --- Countering the Swarm: Protecting the Joint Force in
  the Drone Age (Sep 2025)\\
  \url{https://s3.us-east-1.amazonaws.com/files.cnas.org/documents/Report_CUAS_Defense_Sep-2025_final.pdf}
\item
  \textbf{CNA} --- PRC Concepts for UAV Swarms in Future Warfare (Oct
  2025)\\
  \url{https://www.cna.org/analyses/2025/10/prc-concepts-for-uav-swarms-in-future-warfare}
\item
  \textbf{U.S. Secretary of Defense Memorandum} --- Replicator 2
  Direction and Execution (Sep 27, 2024)\\
  \url{https://s3.us-gov-west-1.amazonaws.com/assets.diu.mil/3nanhbfkr0pc/1dkJGhMeAgPldz1nnIwabK/abf85531a4281cddab6b0d8c953440e2/REPLICATOR-2-MEMO-SD-SIGNED__1_.pdf}
\item
  \textbf{X source pointer (The Information)} --- post link queued for
  authenticated extraction\\
  \url{https://x.com/theinformation/status/2021707053727662407?s=46}
\end{enumerate}

\hypertarget{academic-source-library-topic-grouped}{%
\subsection{14) Academic source library
(topic-grouped)}\label{academic-source-library-topic-grouped}}

\hypertarget{marl-core}{%
\subsubsection{MARL core}\label{marl-core}}

\begin{enumerate}
\def\labelenumi{\arabic{enumi}.}
\tightlist
\item
  Multi-Agent Reinforcement Learning: Independent vs.~Cooperative Agents
  (1993), cited by 1789.
  \url{https://doi.org/10.1016/b978-1-55860-307-3.50049-6}
\item
  Counterfactual Multi-Agent Policy Gradients (2018), cited by 1543.
  \url{https://doi.org/10.1609/aaai.v32i1.11794}
\item
  Multi-Agent Actor-Critic for Mixed Cooperative-Competitive
  Environments (2017), cited by 1014.
  \url{http://arxiv.org/abs/1706.02275}
\item
  Cooperative Multi-agent Control Using Deep Reinforcement Learning
  (2017), cited by 846.
  \url{https://doi.org/10.1007/978-3-319-71682-4_5}
\item
  Multi-agent deep reinforcement learning: a survey (2021), cited by
  726. \url{https://doi.org/10.1007/s10462-021-09996-w}
\item
  QMIX: Monotonic Value Function Factorisation for Deep Multi-Agent
  Reinforcement Learning (2018), cited by 475.
  \url{http://arxiv.org/abs/1803.11485}
\item
  Monotonic Value Function Factorisation for Deep Multi-Agent
  Reinforcement Learning (2020), cited by 429.
  \url{http://arxiv.org/abs/2003.08839}
\item
  A review of cooperative multi-agent deep reinforcement learning
  (2022), cited by 404. \url{https://doi.org/10.1007/s10489-022-04105-y}
\item
  Multi-Agent Reinforcement Learning: A Review of Challenges and
  Applications (2021), cited by 315.
  \url{https://doi.org/10.3390/app11114948}
\item
  Robust Multi-Agent Reinforcement Learning via Minimax Deep
  Deterministic Policy Gradient (2019), cited by 279.
  \url{https://doi.org/10.1609/aaai.v33i01.33014213}
\item
  QTRAN: Learning to Factorize with Transformation for Cooperative
  Multi-Agent Reinforcement Learning (2019), cited by 222.
  \url{https://doi.org/10.48550/arxiv.1905.05408}
\item
  Multi-Agent Deep Reinforcement Learning for Multi-Robot Applications:
  A Survey (2023), cited by 140. \url{https://doi.org/10.3390/s23073625}
\item
  Dealing with Non-Stationarity in Multi-Agent Deep Reinforcement
  Learning (2019), cited by 124. \url{http://arxiv.org/abs/1906.04737}
\item
  Value-Decomposition Networks For Cooperative Multi-Agent Learning
  (2017). \url{https://arxiv.org/abs/1706.05296}
\item
  Learning to Communicate with Deep Multi-Agent Reinforcement Learning
  (2016). \url{https://arxiv.org/abs/1605.06676}
\end{enumerate}

\hypertarget{swarm-robotics-human-swarm-interaction}{%
\subsubsection{Swarm robotics + human-swarm
interaction}\label{swarm-robotics-human-swarm-interaction}}

\begin{enumerate}
\def\labelenumi{\arabic{enumi}.}
\tightlist
\item
  Swarm robotics: a review from the swarm engineering perspective
  (2013), cited by 1638. \url{https://doi.org/10.1007/s11721-012-0075-2}
\item
  A Survey on Aerial Swarm Robotics (2018), cited by 604.
  \url{https://doi.org/10.1109/tro.2018.2857475}
\item
  Swarm of micro flying robots in the wild (2022), cited by 490.
  \url{https://doi.org/10.1126/scirobotics.abm5954}
\item
  A review of swarm robotics tasks (2015), cited by 428.
  \url{https://doi.org/10.1016/j.neucom.2015.05.116}
\item
  Swarm Robotic Behaviors and Current Applications (2020), cited by 413.
  \url{https://doi.org/10.3389/frobt.2020.00036}
\item
  Human Interaction With Robot Swarms: A Survey (2015), cited by 387.
  \url{https://doi.org/10.1109/thms.2015.2480801}
\item
  Swarm Robotics: Past, Present, and Future {[}Point of View{]} (2021),
  cited by 327. \url{https://doi.org/10.1109/jproc.2021.3072740}
\item
  Search and tracking algorithms for swarms of robots: A survey (2015),
  cited by 180. \url{https://doi.org/10.1016/j.robot.2015.08.010}
\item
  Predictive control of aerial swarms in cluttered environments (2021),
  cited by 151. \url{https://doi.org/10.1038/s42256-021-00341-y}
\item
  Human swarm interaction for radiation source search and localization
  (2008), cited by 112. \url{https://doi.org/10.1109/sis.2008.4668287}
\item
  Human Swarm Interaction: An Experimental Study of Two Types of
  Interaction with Foraging Swarms (2013), cited by 104.
  \url{https://doi.org/10.5898/jhri.2.2.kolling}
\item
  Mean-field models in swarm robotics: a survey (2019), cited by 94.
  \url{https://doi.org/10.1088/1748-3190/ab49a4}
\end{enumerate}

\hypertarget{llm-agents-orchestration}{%
\subsubsection{LLM agents +
orchestration}\label{llm-agents-orchestration}}

\begin{enumerate}
\def\labelenumi{\arabic{enumi}.}
\tightlist
\item
  A survey on large language model based autonomous agents (2024), cited
  by 797. \url{https://doi.org/10.1007/s11704-024-40231-1}
\item
  A survey on LLM-based multi-agent systems: workflow, infrastructure,
  and challenges (2024), cited by 153.
  \url{https://doi.org/10.1007/s44336-024-00009-2}
\item
  AutoGen: Enabling Next-Gen LLM Applications via Multi-Agent
  Conversation (2023), cited by 138.
  \url{http://arxiv.org/abs/2308.08155}
\item
  ChatEval: Towards Better LLM-based Evaluators through Multi-Agent
  Debate (2023), cited by 51. \url{http://arxiv.org/abs/2308.07201}
\item
  Large Language Model based Multi-Agents: A Survey of Progress and
  Challenges (2024), cited by 50. \url{http://arxiv.org/abs/2402.01680}
\item
  Exploring Large Language Model based Intelligent Agents: Definitions,
  Methods, and Prospects (2024), cited by 35.
  \url{http://arxiv.org/abs/2401.03428}
\item
  Personal LLM Agents: Insights and Survey about the Capability,
  Efficiency and Security (2024), cited by 30.
  \url{http://arxiv.org/abs/2401.05459}
\item
  Understanding the planning of LLM agents: A survey (2024), cited by
  26. \url{http://arxiv.org/abs/2402.02716}
\item
  Encouraging Divergent Thinking in Large Language Models through
  Multi-Agent Debate (2023), cited by 26.
  \url{http://arxiv.org/abs/2305.19118}
\item
  LLM-Based Multi-Agent Systems for Software Engineering: Literature
  Review, Vision and the Road Ahead (2024), cited by 8.
  \url{http://arxiv.org/abs/2404.04834}
\item
  CAMEL: Communicative Agents for Mind Exploration of Large Language
  Model Society (2023). \url{https://arxiv.org/abs/2303.17760}
\item
  AgentVerse: Facilitating Multi-Agent Collaboration and Exploring
  Emergent Behaviors (2023). \url{https://arxiv.org/abs/2308.10848}
\item
  MetaGPT: Meta Programming for Multi-Agent Collaborative Framework
  (2023). \url{https://arxiv.org/abs/2308.00352}
\item
  SWE-agent: Agent-Computer Interfaces Enable Automated Software
  Engineering (2024). \url{https://arxiv.org/abs/2405.15793}
\end{enumerate}

\hypertarget{safety-governance-and-legal-accountability}{%
\subsubsection{Safety, governance, and legal
accountability}\label{safety-governance-and-legal-accountability}}

\begin{enumerate}
\def\labelenumi{\arabic{enumi}.}
\tightlist
\item
  A comprehensive survey on safe reinforcement learning (2015), cited by
  1191.
  \url{https://jmlr.csail.mit.edu/papers/volume16/garcia15a/garcia15a.pdf}
\item
  Safe Learning in Robotics: From Learning-Based Control to Safe
  Reinforcement Learning (2022), cited by 594.
  \url{https://doi.org/10.1146/annurev-control-042920-020211}
\item
  Trustworthy artificial intelligence (2020), cited by 507.
  \url{https://doi.org/10.1007/s12525-020-00441-4}
\item
  Meaningful Human Control over Autonomous Systems: A Philosophical
  Account (2018), cited by 405.
  \url{https://doi.org/10.3389/frobt.2018.00015}
\item
  On banning autonomous weapon systems: human rights, automation, and
  the dehumanization of lethal decision-making (2012), cited by 339.
  \url{https://doi.org/10.1017/s1816383112000768}
\item
  A Review of Trustworthy and Explainable Artificial Intelligence (XAI)
  (2023), cited by 225.
  \url{https://doi.org/10.1109/access.2023.3294569}
\item
  Autonomous weapons systems, killer robots and human dignity (2018),
  cited by 176. \url{https://doi.org/10.1007/s10676-018-9494-0}
\item
  Autonomous Weapons and International Humanitarian Law: Advantages,
  Open Technical Questions and Legal Issues to be Clarified (2014),
  cited by 113. \url{https://archive-ouverte.unige.ch/unige:37976}
\item
  Autonomous Weapon Systems and International Humanitarian Law: A Reply
  to the Critics (2012), cited by 101.
  \url{https://doi.org/10.2139/ssrn.2184826}
\item
  Open Problems and Fundamental Limitations of Reinforcement Learning
  from Human Feedback (2023), cited by 88.
  \url{http://arxiv.org/abs/2307.15217}
\end{enumerate}

\hypertarget{adversarial-robustness-and-resilience}{%
\subsubsection{Adversarial robustness and
resilience}\label{adversarial-robustness-and-resilience}}

\begin{enumerate}
\def\labelenumi{\arabic{enumi}.}
\tightlist
\item
  Secure impulsive synchronization control of multi-agent systems under
  deception attacks (2018), cited by 324.
  \url{https://doi.org/10.1016/j.ins.2018.04.020}
\item
  Secure impulsive synchronization in Lipschitz-type multi-agent systems
  subject to deception attacks (2020), cited by 165.
  \url{https://doi.org/10.1109/jas.2020.1003297}
\item
  A study on cyber-security of autonomous and unmanned vehicles (2015),
  cited by 136. \url{https://doi.org/10.1177/1548512915575803}
\item
  Cooperative adaptive fault-tolerant control for multi-agent systems
  with deception attacks (2020), cited by 111.
  \url{https://doi.org/10.1016/j.jfranklin.2019.12.032}
\item
  Detection of Cyber-attacks to indoor real time localization systems
  for autonomous robots (2017), cited by 87.
  \url{https://doi.org/10.1016/j.robot.2017.10.006}
\item
  Fault-Tolerant secure consensus tracking of delayed nonlinear
  multi-agent systems with deception attacks and uncertain parameters
  via impulsive control (2019), cited by 85.
  \url{https://doi.org/10.1016/j.cnsns.2019.105043}
\item
  Enhancing Autonomous System Security and Resilience With Generative
  AI: A Comprehensive Survey (2024), cited by 60.
  \url{https://doi.org/10.1109/access.2024.3439363}
\item
  Failure-Scenario Maker for Rule-Based Agent using Multi-agent
  Adversarial Reinforcement Learning and its Application to Autonomous
  Driving (2019), cited by 57.
  \url{https://doi.org/10.24963/ijcai.2019/832}
\item
  Multi-Agent Adversarial Inverse Reinforcement Learning (2019), cited
  by 45. \url{http://arxiv.org/abs/1907.13220}
\item
  On the Robustness of Cooperative Multi-Agent Reinforcement Learning
  (2020). \url{https://arxiv.org/abs/2006.07538}
\end{enumerate}

\hypertarget{build-target-for-review}{%
\subsection{Build target for review}\label{build-target-for-review}}

\begin{itemize}
\tightlist
\item
  Generated \texttt{.tex} source path (internal):
  \texttt{Whitepaper/writing/notes/tex/2026-02-12-annex-a-supporting-notes-v6.tex}
\item
  Generated PDF artifact path (internal):
  \texttt{out/memos/2026-02-12-annex-a-supporting-notes-v6.pdf}
\item
  Publication status: \textbf{draft / non-published}.
\end{itemize}

\end{document}
